\documentclass[12pt]{article}
\usepackage[margin=0.58in,top=0.62in,bottom=0.48in]{geometry}
\usepackage{lmodern}
\usepackage{microtype}
\usepackage{fancyhdr}
\usepackage{titlesec}
\usepackage{enumitem}
\usepackage{lastpage}
\usepackage[hidelinks]{hyperref}

\setlength{\parindent}{0pt}
\setlength{\parskip}{0.18em}
\setlength{\headheight}{26pt}
\raggedbottom
\titleformat{\section}{\normalsize\bfseries\scshape}{}{0pt}{}
\titleformat{\subsection}{\normalsize\bfseries}{}{0pt}{}
\titlespacing*{\section}{0pt}{0.35em}{0.12em}
\titlespacing*{\subsection}{0pt}{0.25em}{0.08em}
\pagestyle{fancy}
\fancyhf{}
\fancyhead[C]{\footnotesize Lit Example Reader \quad Assignment ID: U1L13LE \quad Page \thepage\ of \pageref{LastPage}}
\renewcommand{\headrulewidth}{0.5pt}

\newenvironment{playpassage}{\begin{quote}\footnotesize\setlength{\parskip}{0.08em}\setlength{\parindent}{0pt}}{\end{quote}}
\newcommand{\speaker}[1]{\textbf{#1.}}

\begin{document}

\begin{center}
{\LARGE\bfseries Week 13 Lit Example Reader}\par\vspace{0.02em}
{\large\scshape Shakespeare: Testimony, Sources, and Credibility}\par\vspace{0.06em}
\rule{0.78\textwidth}{0.5pt}
\end{center}

\textbf{Purpose:} Shakespeare passages for Week 13: judge testimony by access, expertise, bias, corroboration, and careful conclusion.

\textbf{What to watch for:} Who is speaking, what the speaker could know, what pressure shapes the testimony, and whether other evidence confirms it.

\section*{Before the Passages: The Week 13 Logic Tool}

\begin{itemize}[leftmargin=1.3em,itemsep=0.05em,topsep=0.08em,parsep=0.03em]
\item \textbf{Source:} the person, record, or voice giving testimony.
\item \textbf{Access:} whether the source was in a position to know.
\item \textbf{Bias or pressure:} a reason the source may distort or steer the hearer.
\item \textbf{Corroboration:} independent support or conflict from other evidence.
\end{itemize}

\section*{Passage 1: Macbeth, Act 1, Scene 3 --- The Witches as Sources}

\textbf{Context:} Macbeth and Banquo hear the witches predict titles and future greatness. Part of the testimony is soon confirmed, but the source remains dangerous and unclear.

\begin{playpassage}
\speaker{First Witch} All hail, Macbeth! hail to thee, Thane of Glamis!\\
\speaker{Second Witch} All hail, Macbeth! hail to thee, Thane of Cawdor!\\
\speaker{Third Witch} All hail, Macbeth, that shalt be king hereafter!\\
...\\
\speaker{Banquo} If you can look into the seeds of time,\\
And say which grain will grow and which will not,\\
Speak then to me.\\
...\\
\speaker{Ross} And, for an earnest of a greater honour,\\
He bade me, from him, call thee Thane of Cawdor.
\end{playpassage}

\subsection*{Reader Notes}

This is your assisted example. The witches have surprising access to something Macbeth did not know: he has been named Thane of Cawdor. That corroboration makes their testimony tempting. But credibility is not the same as one accurate detail. Their source of knowledge is unknown, their wording is partial, and their moral pressure pushes Macbeth toward ambition. A careful conclusion might be: ``The Cawdor claim is corroborated; the kingship claim is not yet proved, and the source remains dangerous.''

\textbf{Reading task:} Mark the corroborated claim and the uncorroborated claim. In the margin, write one reason the witches are tempting sources and one reason they remain unreliable.

\newpage

\section*{Passage 2: Julius Caesar, Act 3, Scene 2 --- Brutus Testifies About Caesar}

\textbf{Context:} After Caesar's murder, Brutus tells the crowd why Caesar deserved death. The crowd must decide how much to trust Brutus's testimony about Caesar's ambition.

\begin{playpassage}
\speaker{Brutus} If there be any in this assembly, any dear friend of Caesar's, to him I say, that Brutus' love to Caesar was no less than his. If then that friend demand why Brutus rose against Caesar, this is my answer: Not that I loved Caesar less, but that I loved Rome more. Had you rather Caesar were living, and die all slaves, than that Caesar were dead, to live all free men? As Caesar loved me, I weep for him; as he was fortunate, I rejoice at it; as he was valiant, I honour him; but, as he was ambitious, I slew him.
\end{playpassage}

\subsection*{Reader Notes}

This passage shifts more work to you. Brutus has relevant access because he knew Caesar and took part in public life. He may also be sincere. But he has strong pressure to justify an assassination he helped commit. Mark what Brutus actually testifies to, what evidence he gives in the speech, and what would corroborate or challenge his claim.

\textbf{Reading task:} Underline Brutus's central claim about Caesar. Next to the passage, write one strength and one weakness in Brutus's credibility. Save the final trust judgment for the worksheet.

\section*{How to Use These Passages on the Worksheet}

Do not summarize the plot. Use Week 13's credibility map: claim, source, access, expertise or competence, bias, corroboration, and careful conclusion.

\end{document}
